\title{ssrn 5333703}
\begin{document}
\maketitle

\begin{abstract}
This paper develops a dynamic general equilibrium model to analyze the macroe- conomic and distributive consequences of artificial general intelligence (AGI) diffusion. Building on task-based and sectorally heterogeneous production, we model AGI capi- tal as a general-purpose technology that substitutes asymmetrically for human labor across sectors and skill classes. The model incorporates endogenous AGI investment, nonlinear diffusion dynamics, and labor displacement, along with a fiscal mechanism that taxes AGI-generated value added and redistributes it via targeted universal basic income (UBI). We show that AGI adoption leads to a secular decline in labor income and a rising concentration of capital rents, generating structural inequality. Absent intervention, the economy converges to a post-labor steady state marked by vanishing wages and increasing income concentration. However, an AGI-linked value-added tax, rebated to displaced workers through targeted UBI, can fully offset income losses and improve intertemporal welfare. Despite transitional frictions due to capital adjustment costs, the model admits a balanced growth path with sustained output growth in the absence of human labor. These results provide a tractable framework for evaluating automation-era fiscal architecture and offer policy-relevant insights into redistribution, labor displacement, and structural macroeconomic adjustment.
\end{abstract}

\section{Recovered Text}
Artificial General Intelligence and the Sectoral
Transition to a Post-Labor Economy: A Dynamic
General Equilibrium Analysis
Pascal Stiefenhofer
Newcastle Univerity, United Kingdom
Department of Economics
June 25, 2025
1
This preprint research paper has not been peer reviewed. Electronic copy available at: https://ssrn.com/abstract=5333703
Preprint not peer reviewed
Artificial General Intelligence and the Sectoral
Transition to a Post-Labor Economy: A Dynamic
General Equilibrium Analysis
Abstract
This paper develops a dynamic general equilibrium model to analyze the macroe-
conomic and distributive consequences of artificial general intelligence (AGI) diffusion.
Building on task-based and sectorally heterogeneous production, we model AGI capi-
tal as a general-purpose technology that substitutes asymmetrically for human labor
across sectors and skill classes. The model incorporates endogenous AGI investment,
nonlinear diffusion dynamics, and labor displacement, along with a fiscal mechanism
that taxes AGI-generated value added and redistributes it via targeted universal basic
income (UBI). We show that AGI adoption leads to a secular decline in labor income
and a rising concentration of capital rents, generating structural inequality. Absent
intervention, the economy converges to a post-labor steady state marked by vanishing
wages and increasing income concentration. However, an AGI-linked value-added tax,
rebated to displaced workers through targeted UBI, can fully offset income losses and
improve intertemporal welfare. Despite transitional frictions due to capital adjustment
costs, the model admits a balanced growth path with sustained output growth in the
absence of human labor. These results provide a tractable framework for evaluating
automation-era fiscal architecture and offer policy-relevant insights into redistribution,
labor displacement, and structural macroeconomic adjustment.
keywords:
Artificial General Intelligence (AGI); Automation; Universal Basic Income
(UBI); Value-Added Tax (VAT); General Equilibrium; Labor Displacement
1
Introduction
The accelerating progress in artificial general intelligence (AGI) promises to redefine the
structure of production, the nature of labor, and the foundations of economic distribution.
Unlike task-specific automation or narrow AI, AGI entails the development of broadly ca-
pable, adaptable systems that can substitute for human labor across a wide spectrum of
cognitive and manual tasks (Russell and Norvig, 2010; Goertzel, 2007; Brynjolfsson and
1
This preprint research paper has not been peer reviewed. Electronic copy available at: https://ssrn.com/abstract=5333703
Preprint not peer reviewed
McAfee, 2014; Korinek, 2022). The implications of this shift for employment, growth, and
inequality are profound and demand new frameworks that go beyond conventional modeling
of labor-displacing technologies.
A large body of work has analyzed the economic impact of automation. Pioneering task-
based models by Acemoglu and Restrepo (2018, 2020, 2022) demonstrate that automation
can reduce labor demand and shift the labor income share while boosting productivity.
Empirical studies such as Autor (2015), Graetz and Michaels (2018), and Bloom et al.
(2020) show that automation contributes to labor market polarization and rising inequality,
with significant sectoral heterogeneity. However, these studies typically treat automation as
exogenous or partial, and do not address the economy-wide implications of general-purpose
AGI technologies.
In contrast to earlier work emphasizing task-level automation, recent macroeconomic re-
search has begun to engage directly with the broader implications of artificial general intel-
ligence (AGI). Korinek (2021), Korinek and Stiglitz (2021), and Stiefenhofer (2025) explore
how AGI may induce large-scale technological unemployment, raising profound distribu-
tional and institutional challenges. Trajtenberg (2018) characterizes AI as a general-purpose
technology (GPT) with wide-ranging spillover effects across sectors, while Brynjolfsson et al.
(2022) emphasize its potential to raise productivity through complementary innovations.
Models of endogenous automation, such as Guerreiro et al. (2020), formalize firms’ technol-
ogy adoption decisions as equilibrium responses to relative prices and factor substitution.
These approaches illuminate the general equilibrium impact of automation but typically ab-
stract from redistributive policy design and sectoral heterogeneity. Stiefenhofer and Chen
(2024) extend this literature by analyzing the dynamic stability of industrial artificial intel-
ligence within Cobb–Douglas production structures, offering early insight into the long-run
macroeconomic behavior of AGI-driven economies.
Several new contributions further highlight the urgency of analyzing AGI’s systemic im-
pact. OECD (2024) and IMF (2024) provide institutional assessments of AI’s macroeconomic
and distributional effects. Acemoglu (2024) offers a stylized macroeconomic framework cap-
turing key substitution and complementarity forces driven by AGI. Windmann et al. (2024)
and Lee and Su (2025) explore how industrial AI challenges integration into existing systems,
while Erdil et al. (2025) propose a novel dynamic model (GATE) linking AI, growth, and
redistribution. Subaveerapandiyan and Shimray (2024) analyze the bibliometric evolution
of AI-labor studies, emphasizing the field’s expanding scope.
In the context of general equilibrium modeling, automation has been integrated into
heterogeneous agent frameworks such as Krusell and Smith (2000), who analyze inequality
dynamics, and Hubmer et al. (2020), who study the role of capital taxation and wealth
concentration.
Jones (2016) and Guner et al. (2023) extend endogenous growth models
to examine the macroeconomic effects of AI and automation.
A related strand of work
focuses on the design of redistributive institutions under automation, including proposals
for universal basic income (UBI) (Bidadanure, 2019; Hoynes and Rothstein, 2019; Gentilini,
2020).
Werner (2021) and Lowrey (2018) explore how UBI interacts with technological
displacement and labor market participation. However, there is limited work linking AGI-
induced productivity gains to financing mechanisms for redistribution. Notable exceptions
include Koenig et al. (2022), who incorporate optimal transfers in a GE setting, and Korinek
(2021), who model the fiscal implications of taxing AI-generated surplus.
2
This preprint research paper has not been peer reviewed. Electronic copy available at: https://ssrn.com/abstract=5333703
Preprint not peer reviewed
To advance this literature, we adopt a dynamic, sectorally heterogeneous general equi-
librium framework in which AGI capital is both a productive and taxable input. A general
equilibrium model is especially suited to studying AGI because it captures how labor dis-
placement, capital reallocation, and policy interventions—such as value-added taxation and
universal basic income—interact across markets and sectors simultaneously (Daruich and
Fern´andez, 2021; Imrohoroglu et al., 2021; Nir et al., 2021). Unlike partial equilibrium ap-
proaches, GE analysis allows for endogenous feedbacks in factor prices, income distribution,
and fiscal flows, which are essential when technological change is pervasive and nonlinear
(Jones, 2021; Lofgren et al., 2002). For example, Daruich and Fern´andez (2021) analyze
a UBI-financed transition in a dynamic GE model with automation shocks, showing that
transition dynamics critically affect intergenerational welfare. Similarly, Imrohoroglu et al.
(2021) and Nir et al. (2021) use GE frameworks to explore how automation and policy
affect sectoral labor displacement, capital accumulation, and occupational choice. These
studies demonstrate that a general equilibrium setting is indispensable for capturing both
the short-run disruption and long-run restructuring induced by AGI. In particular, the GE
structure accommodates sector-specific automation dynamics, general-purpose technology
spillovers, and the reallocation of displaced labor and capital—making it a robust platform
for analyzing both transitional and steady-state outcomes under AGI diffusion.
AGI diffusion is modeled via sector-specific logistic functions, generating time-varying
exposure to automation across labor classes. Output is produced through a nested CES
aggregator combining human labor, AGI labor (modeled as capital-embedded services), and
two forms of capital.
We introduce a value-added tax (VAT) on AGI-attributed output
and use the proceeds to finance a targeted universal basic income calibrated to displaced
labor. Importantly, we incorporate nonlinear AGI capital adjustment costs and derive tran-
sitional dynamics consistent with empirical observations of technological diffusion (Comin
and Hobijn, 2010; Bresnahan and Trajtenberg, 1995).
Our contribution builds on task-based frameworks (Acemoglu and Restrepo, 2019; Re-
strepo, 2022), general equilibrium analyses of inequality (Krusell and Smith, 2000; Hubmer
et al., 2020), and macro-fiscal redistribution models (Korinek, 2021; Aksoy et al., 2023). In
contrast to prior work that treats AGI as a static productivity shift or exogenous technology,
we explicitly embed AGI in the structure of production, finance, and income distribution.
We formally derive conditions under which AGI-VAT revenues can sustain income replace-
ment for displaced labor, characterize the emergence of balanced growth without human
labor, and highlight transitional disequilibria induced by capital adjustment frictions.
By embedding AGI into both the production and fiscal architecture of a dynamic econ-
omy, the model provides a platform for analyzing long-run welfare outcomes, inequality
trajectories, and the fiscal viability of universal transfers. It also raises normative questions
about the role of the state in governing AGI diffusion and mediating its gains, particularly
in economies where income becomes increasingly decoupled from labor (Piketty, 2014; Farhi
and Werning, 2019).
This paper is organized as follows. Section 2 introduces the dynamic general equilibrium
model featuring AGI diffusion, value-added taxation (VAT), and targeted universal basic
income (UBI). Section 3 presents the analytical results, focusing on automation dynamics,
distributional shifts, and transitional general equilibrium properties. Section 4 discusses the
policy implications and situates the findings within the existing academic literature. Section
3
This preprint research paper has not been peer reviewed. Electronic copy available at: https://ssrn.com/abstract=5333703
Preprint not peer reviewed
5 concludes. Technical results and some proofs are provided in the Appendix.
2
Model Overview
We consider a continuous-time economy over t ∈[0, ∞), populated by agents partitioned
into labor classes j ∈\{c, m, h\} and economic sectors i = 1, . . . , n. The total population is
given by
N =
X
j∈\{c,m,h\}
Nj,
(1)
where Nj denotes the size of labor class j: cognitive workers (j = c, early automation risk),
manual workers (j = m, late automation risk), and hybrid workers (j = h, intermediate
risk). Labor is distributed across sectors according to the matrix β = [βij], where βij ≥0 is
the fraction of labor class j employed in sector i, satisfying
X
i
βij = 1,
∀j.
(2)
Automation intensity in sector i follows a logistic diffusion process
αi(t) =
1
1 + e−κi(t−t∗
i ),
(3)
where κi > 0 is the sector-specific diffusion rate and t∗
i is the inflection point. This functional
form is supported by empirical studies on the adoption of general-purpose technologies such
as electricity, IT, and AI (Griliches, 1957; Mansfield, 1961; Comin and Hobijn, 2010). It
captures slow initial adoption, rapid acceleration, and eventual saturation. The automation
exposure of labor class j is a weighted average of sectoral automation intensities
γj(t) =
X
i
βijαi(t).
(4)
The effective labor supply of class j in sector i is then
Lij(t) = βijNj
1 −γj(t)

,
(5)
reflecting that γj(t) reduces the productive share of labor due to automation. Total effective
human labor in sector i is
Li,h(t) =
X
j
Lij(t).
(6)
AGI labor input in sector i is
Li,AGI(t) = ψiαi(t)K(AGI)
i
(t),
(7)
where ψi > 0 is a productivity parameter, K(AGI)
i
(t) is AGI-specific capital, and αi(t) de-
termines its activation level. This expression defines AGI labor not as a traditional labor
input, but as a derived, capital-embedded service flow that substitutes for human labor.
4
This preprint research paper has not been peer reviewed. Electronic copy available at: https://ssrn.com/abstract=5333703
Preprint not peer reviewed
Specifically, K(AGI)
i
(t) represents the stock of AGI-enabled capital (e.g., autonomous sys-
tems, intelligent software), and αi(t) denotes the sector’s automation intensity—governing
the degree to which this capital is operationalized. The parameter ψi > 0 captures task-
specific or sectoral productivity differences in translating AGI capital into effective labor
services. The structure reflects the idea that AGI labor is activated capital: without au-
tomation infrastructure (αi(t)), AGI capital is idle; but as automation spreads, AGI capital
increasingly delivers labor-equivalent functionality. This approach aligns with task-based
models of automation (Acemoglu and Restrepo, 2019) and extends capital–labor substitu-
tion by introducing endogenous, technology-driven activation of digital labor inputs. Unlike
traditional labor, AGI labor is not constrained by population, but by the pace of technolog-
ical diffusion and investment in specialized capital.
Production in sector i is modeled using a three-level nested CES aggregator:
Qi(t) = Ai

θ1 (Li,h(t)ρ1 + Li,AGI(t)ρ1)
η1
ρ1
+ θ2

(K(H)
i
(t))ρ2 + (K(AGI)
i
(t))ρ2 η2
ρ2
 1
η
,
(8)
where Ai > 0 denotes total factor productivity, and θ1+θ2 = 1 are share parameters. The
parameters ρ1 and ρ2 govern the curvature of the respective CES nests for labor and capital
inputs. The inner substitution elasticities are given by 1/(1 −η1) and 1/(1 −η2), capturing
the substitutability between human and AGI labor and between traditional and AGI-specific
capital, respectively. The outer elasticity parameter η regulates the substitutability between
the two composite bundles—labor and capital—allowing for sectoral flexibility in the al-
location of value added between factor types. This three-layered structure captures both
the within-factor substitution effects and the broader reallocation dynamics that character-
ize AGI-driven automation, consistent with empirical CES specifications in the growth and
technology diffusion literature.
The nested CES structure reflects a nuanced representation of production under AGI-
driven automation. The innermost labor nest captures the idea that human labor Li,h(t)
and AGI labor Li,AGI(t) are imperfect substitutes in performing tasks. The substitution
parameter ρ1 governs the elasticity between these inputs, enabling the model to span a
continuum between automation-augmenting environments—where human and AGI labor
are only weakly substitutable—and automation-replacing regimes, in which AGI increasingly
displaces human input (Acemoglu and Restrepo, 2019).
Parallel to this, the capital nest aggregates traditional physical capital K(H)
i
(t) and AGI-
specific capital K(AGI)
i
(t).
This formulation acknowledges the growing role of intelligent
capital assets, such as autonomous systems and machine learning platforms, in production.
The elasticity parameter ρ2 determines the degree of substitution between conventional and
AGI-embedded capital, aligning with frameworks in which capital heterogeneity responds
endogenously to technological progress (Acemoglu and Restrepo, 2022).
At the outer level, the model combines the labor and capital composites through a CES
aggregator governed by the elasticity parameter η. This outer nest allows for variation in
how sectors substitute between total labor-type and capital-type inputs, accommodating
structural asymmetries and sectoral diversity in the adoption of automation. The form is
5
This preprint research paper has not been peer reviewed. Electronic copy available at: https://ssrn.com/abstract=5333703
Preprint not peer reviewed
consistent with two-level CES specifications commonly used in empirical studies of produc-
tion and technological change (Krusell and Smith, 2000).
The share parameters θ1 and θ2, which satisfy θ1 + θ2 = 1, determine the relative weight
of the labor and capital nests in sectoral value added. These weights capture differences in
production composition across sectors, such as between labor-intensive services and capital-
intensive manufacturing, and facilitate comparative statics as automation and capital accu-
mulation evolve over time.
This production architecture is particularly well-suited for modeling AGI as a labor-
equivalent input that is nonetheless capital-embodied.
By explicitly nesting AGI labor
alongside human labor and pairing AGI-specific capital with traditional capital, the model
captures the dual character of intelligent systems: as both cognitive substitutes for human
labor and as physical assets that require capital investment.
This design integrates the
model with task-based theories of automation and endogenous technological change, while
remaining tractable within calibrated macroeconomic frameworks.
In each sector i, output prices are determined under the standard assumption of perfect
competition. Firms behave as price takers and set prices equal to marginal cost. Let Ci(Qi(t))
denote the total cost function associated with producing quantity Qi(t), which accounts for
expenditures on effective labor, traditional capital, AGI-specific capital, and any sector-
specific adjustment costs. The marginal cost is then defined as the derivative of the total
cost with respect to output:
MCi(t) = ∂Ci(Qi(t))
∂Qi(t)
.
(9)
Under the zero-profit condition characteristic of competitive equilibrium, firms set prices
equal to marginal cost:
Pi(t) = MCi(t) = ∂Ci(Qi(t))
∂Qi(t)
.
(10)
This pricing rule follows directly from microeconomic duality theory. Specifically, it can
be derived from the firm’s cost minimization problem via Shepard’s Lemma
∂Ci(Qi(t), p(t))
∂Qi(t)
= λ(t),
(11)
where λ(t) is the Lagrange multiplier on the production constraint and represents the
marginal cost of output. The input price vector p(t) includes the cost of both traditional and
AGI-related inputs, making this formulation sensitive to changes in automation and tech-
nological intensity. In the nested CES framework introduced in Equation (8), the marginal
cost function is not available in closed form due to the nonlinear interactions among labor
types, capital types, and their substitution elasticities. Instead, MCi(t) is implicitly de-
fined and varies with the sectoral automation intensity, the degree of factor substitutability,
and the evolving composition of capital. Accordingly, numerical methods—such as local
log-linearization or calibrated simulations—are typically employed to characterize price dy-
namics and analyze equilibrium responses to policy or technological shocks.
6
This preprint research paper has not been peer reviewed. Electronic copy available at: https://ssrn.com/abstract=5333703
Preprint not peer reviewed
Capital accumulation in each sector i is governed by the following differential equations
˙K(H)
i
(t) = I(H)
i
(t) −δ(H)
i
K(H)
i
(t),
(12)
˙K(AGI)
i
(t) = I(AGI)
i
(t) −ϕi

I(AGI)
i
(t), K(AGI)
i
(t)

,
(13)
where K(H)
i
(t) and K(AGI)
i
(t) denote the sector-specific stocks of traditional and AGI-related
capital, respectively, and I(H)
i
(t), I(AGI)
i
(t) are the corresponding gross investments. The
parameter δ(H)
i
∈(0, 1) is the constant depreciation rate for traditional capital. In contrast,
the evolution of AGI capital reflects nonlinear adjustment dynamics. Depreciation in this
case is modeled as a convex function of both investment and capital
ϕi(I, K) = ϑi
2
 I
K
2
K = ϑi
2
I2
K ,
(14)
where ϑi > 0 is a sector-specific adjustment cost parameter. This functional form implies
that the effective loss of AGI capital increases disproportionately with the rate of investment
relative to the existing capital stock. To see this formally, note that for fixed K > 0, the
depreciation function ϕi(I, K) is strictly convex in I, as the second derivative satisfies
∂2ϕi(I, K)
∂I2
= ϑi
K > 0.
This convexity captures the frictions inherent in AGI deployment—such as organizational re-
structuring, compatibility constraints, training requirements, and coordination delays—which
make rapid scaling of AGI capital increasingly costly. These dynamics are consistent with
the adjustment cost literature on general-purpose technologies (Hayashi, 1982; Acemoglu
and Restrepo, 2022). The use of a nonlinear depreciation function implies that investment
decisions must account not only for marginal capital returns but also for endogenous capi-
tal losses due to scaling inefficiencies. This enriches the model’s dynamics and allows it to
capture realistic transitional frictions in economies undergoing structural automation.
Households in labor class j derive utility from consumption and face income fluctuations
due to sector-specific automation. Lifetime utility is given by the discounted logarithmic
functional form
Uj =
Z ∞
0
e−ρt log
Cj(t)

dt,
(15)
where Cj(t) denotes consumption at time t, and ρ > 0 is the constant rate of time preference.
Households maximize utility subject to a period-by-period budget constraint
Cj(t) ≤
1 −γj(t)

wj(t) + UBI(t),
(16)
where γj(t) ∈[0, 1] represents the automation exposure of labor class j, wj(t) is the class-
specific wage rate, and UBI(t) is a universal basic income transfer financed through taxation
on automated production. The term (1−γj(t)) reflects the employed share of the population
within class j, with higher γj(t) indicating greater technological displacement. Aggregated
consumption across class j is given by
Cagg
j
(t) = Nj
1 −γj(t)

wj(t) + UBI(t)

,
(17)
7
This preprint research paper has not been peer reviewed. Electronic copy available at: https://ssrn.com/abstract=5333703
Preprint not peer reviewed
where Nj is the total mass of individuals in labor class j. For tractability, we assume that all
physical and AGI-specific capital is owned by firms. Households do not earn capital income;
their total income derives exclusively from wages and UBI transfers. This assumption isolates
the distributional effects of automation to the labor market channel and the redistribution
mechanism, abstracting from heterogeneity in asset holdings. Wages are determined by the
marginal productivity of human labor
wj(t) =
X
i
βij
∂Qi(t)
∂Li,h(t),
(18)
where βij denotes the exposure weight of labor class j to sector i, and Qi(t) is the production
function in sector i. This specification implies that a worker’s earnings depend on their
sectoral allocation and the sector-specific marginal product of human labor.
To finance
redistribution, we introduce a value-added tax (VAT) on AGI-enabled production.
The
VAT is designed to target automation-derived surplus, ensuring that firms deploying AGI
contribute in proportion to its productive role. The VAT rate in sector i is specified as
τi(t) = τ0 · Ci,AGI(t)
Pi(t)Qi(t),
(19)
where τ0 ∈[0, 1] is a policy parameter, Ci,AGI(t) denotes the cost of AGI inputs in sector i,
and Pi(t)Qi(t) is the total value of sectoral output. This structure ensures that VAT burdens
scale with the degree of automation, targeting sectors where AGI capital contributes a larger
share of value-added. Total VAT revenue collected from sector i is
Ti(t) = τi(t)Pi(t)Qi(t),
T(t) =
X
i
Ti(t),
(20)
where T(t) is the aggregate VAT revenue available for redistribution. We assume that these
funds are transferred exclusively to displaced workers. The total mass of displaced individuals
is
Ndisp(t) =
X
j
Njγj(t),
(21)
and the UBI transfer per displaced individual is given by
UBI(t) =
(
T(t)
Ndisp(t),
if Ndisp(t) > 0,
0,
otherwise.
(22)
To guarantee the adequacy of this policy, we impose a fiscal constraint requiring that total
VAT revenue must at least replace the wage income lost due to displacement
T(t) ≥
X
j
Njγj(t)wj(t),
(23)
which implies that
UBI(t) ≥
1
Ndisp(t)
X
j
Njγj(t)wj(t).
8
This preprint research paper has not been peer reviewed. Electronic copy available at: https://ssrn.com/abstract=5333703
Preprint not peer reviewed
This ensures that each displaced worker receives a transfer approximating their forgone
wage income, maintaining baseline consumption levels and sustaining aggregate demand
during the automation transition. The mechanism aligns with policy proposals advocating
for automation-adjusted fiscal instruments to manage distributional effects and preserve
macroeconomic stability (Acemoglu and Restrepo, 2022; Korinek, 2021; Guner et al., 2023).
Definition 2.1 (Dynamic General Equilibrium). A dynamic general equilibrium is a collec-
tion of time-dependent functions
\{Cj(t), Pi(t), wj(t), Qi(t), K(H)
i
(t), K(AGI)
i
(t), I(H)
i
(t), I(AGI)
i
(t), Ti(t), UBI(t)\}i,j,t≥0
such that for all t ≥0, the following conditions hold:
(i) Household optimization: For each labor class j,
\{Cj(t)\} = arg max
Cj(·)
Z ∞
0
e−ρt log Cj(t) dt
subject to the present-value intertemporal budget constraint
Z ∞
0
e−ρtCj(t) dt ≤
Z ∞
0
e−ρt [(1 −γj(t))wj(t) + UBI(t)] dt.
(ii) Firm profit maximization: For each sector i, firms choose capital inputs \{K(H)
i
(t), K(AGI)
i
(t)\}
and investment \{I(AGI)
i
(t)\} to maximize
max
\{K(H)
i
(t),K(AGI)
i
(t)\}
Z ∞
0
e−rt

Pi(t)Qi(t) −
X
j
βijNj(1 −γj(t))wj(t)
−rH(t)K(H)
i
(t) −rAGI(t)K(AGI)
i
(t) −ϕi

I(AGI)
i
(t), K(AGI)
i
(t)
 
dt
(24)
subject to the sectoral production function
Qi(t) = Ai

θ1 (Li,h(t)ρ1 + Li,AGI(t)ρ1)
η1
ρ1 + θ2

(K(H)
i
(t))ρ2 + (K(AGI)
i
(t))ρ2 η2
ρ2
 1
η
.
(iii) Capital accumulation:
˙K(H)
i
(t) = I(H)
i
(t) −δ(H)
i
K(H)
i
(t),
˙K(AGI)
i
(t) = I(AGI)
i
(t) −ϕi

I(AGI)
i
(t), K(AGI)
i
(t)

.
(iv) Labor consistency:
Lij(t) = βijNj(1 −γj(t)),
Li,h(t) =
X
j
Lij(t),
Li,AGI(t) = ψiαi(t)K(AGI)
i
(t).
9
This preprint research paper has not been peer reviewed. Electronic copy available at: https://ssrn.com/abstract=5333703
Preprint not peer reviewed
(v) Wage determination:
wj(t) =
X
i
βij
∂Qi(t)
∂Li,h(t).
(vi) Goods market clearing:
X
i
Pi(t)Qi(t) =
X
j
NjCj(t) +
X
i
h
I(H)
i
(t) + I(AGI)
i
(t)
i
.
(vii) Government budget balance:
T(t) = UBI(t) · Ndisp(t),
Ndisp(t) =
X
j
Njγj(t).
(viii) Price setting:
Pi(t) = ∂Ci(Qi(t), p(t))
∂Qi(t)
.
This definition formalizes a dynamic general equilibrium framework designed to capture
the structural transformation of an economy undergoing automation and the integration of
artificial general intelligence (AGI). Households maximize intertemporal utility under income
constraints that reflect declining labor earnings due to technological displacement, partially
offset by redistributive transfers in the form of universal basic income (UBI). The distri-
bution of income across labor classes is shaped by heterogeneous exposure to automation,
introducing dynamic inequality that the fiscal mechanism is designed to mitigate.
Firms make intertemporal investment decisions over traditional and AGI-specific capi-
tal stocks, subject to convex adjustment costs that reflect real-world frictions in deploying
general-purpose technologies. These firms operate in perfectly competitive markets, setting
output prices equal to marginal cost. The production structure incorporates a nested CES
technology that allows for imperfect substitution between human and AGI labor, and be-
tween conventional and AGI-specific capital, thereby capturing the evolving factor intensity
of production as automation deepens.
The government imposes a value-added tax (VAT) on the AGI-attributed share of sectoral
output. Revenues from this tax are recycled into a targeted UBI scheme that compensates
technologically displaced workers. This mechanism operationalizes a core insight from the
literature on automation and fiscal externalities: labor-displacing technologies reduce wage-
based tax bases while concentrating value-added in capital-intensive production, thereby
justifying automation-contingent instruments (Acemoglu and Restrepo, 2022; Broer et al.,
2022; Korinek, 2021). By linking VAT rates to AGI-related cost shares, the model ensures
that more automated sectors contribute proportionally more to redistribution—without di-
rectly taxing capital investment—thus preserving incentives for innovation while internalizing
the distributional costs of automation (Guner et al., 2023; Aksoy et al., 2023).
A fiscal adequacy constraint requires that VAT revenues are sufficient to cover the average
forgone labor income of displaced workers. This condition ensures consumption smoothing
during the transition to a post-labor economy, aligning with normative arguments in the
10
This preprint research paper has not been peer reviewed. Electronic copy available at: https://ssrn.com/abstract=5333703
Preprint not peer reviewed
macro-distributional literature for automation-adjusted transfer schemes (Farhi and Wern-
ing, 2020; Acemoglu and Restrepo, 2019).
Notably, AGI labor is modeled as a capital-
embedded service flow rather than a wage-earning agent, reinforcing the disconnect between
productivity gains and human income in advanced automation regimes and highlighting the
redistributive role of policy.
Overall, the model contributes to the literature on general-purpose technologies (GPTs),
endogenous technological diffusion, and fiscal policy under structural change (Acemoglu
and Restrepo, 2018; Comin and Hobijn, 2010; Krusell and Smith, 2000).
It provides a
theoretically grounded, policy-relevant framework for evaluating the interaction between
automation, taxation, and redistribution under realistic conditions of sectoral heterogeneity
and labor-market asymmetry. By explicitly incorporating an automation-linked VAT and a
displacement-calibrated UBI, the model offers a coherent and tractable basis for designing
equitable transitions in economies confronting AGI-induced disruption.
3
Model Analysis
This section develops key theoretical results that characterize the economic dynamics and
policy implications arising from the diffusion of artificial general intelligence (AGI) across
sectors and labor classes. We analyze how heterogeneous exposure to automation creates
asymmetries in labor displacement and discuss the fiscal mechanisms—such as value-added
taxation on AGI output—available to capture the economic surplus generated by automa-
tion. The analysis further explores the sustainability of redistribution policies like universal
basic income (UBI) funded by AGI-related tax revenues, highlighting their potential to sta-
bilize welfare and inequality in a post-labor economy. Through a series of propositions and
rigorous proofs, we establish fundamental insights into the growth trajectories, transitional
consumption dynamics, and welfare outcomes characterizing the transition to an economy
dominated by AGI capital and labor. These results provide a formal foundation for designing
targeted, class-sensitive policy interventions to manage the disruptive impacts of automation
and ensure equitable economic outcomes.
Proposition 3.1 (Asymmetry in Automation Exposure Across Labor Classes). Let αi(t) =
1
1+e−κi(t−t∗
i ) denote AGI intensity in sector i, and define the automation exposure of labor
class j as
γj(t) =
n
X
i=1
βijαi(t),
(25)
where βij ≥0 and Pn
i=1 βij = 1 for all j. If two labor classes j1 ̸= j2 differ in sectoral exposure
vectors βj1 ̸= βj2, then there exists t > 0 such that γj1(t) ̸= γj2(t). That is, automation
exposure is strictly asymmetric across labor classes with different sectoral weights.
Proof. We restate the definition of AGI intensity in sector i at time t
αi(t) =
1
1 + e−κi(t−t∗
i ),
(26)
11
This preprint research paper has not been peer reviewed. Electronic copy available at: https://ssrn.com/abstract=5333703
Preprint not peer reviewed
a logistic function with parameters κi > 0 and inflection point t∗
i ∈R. This function satisfies
αi(t) ∈(0, 1)
for all t ∈R,
(27)
lim
t→−∞αi(t) = 0,
lim
t→∞αi(t) = 1.
(28)
Each αi(t) is strictly increasing, continuous, and differentiable in t. Automation exposure
for labor class j is defined as a convex combination of sectoral intensities
γj(t) =
n
X
i=1
βijαi(t),
(29)
where βij ≥0 and Pn
i=1 βij = 1 for all j. Let j1 ̸= j2 be two labor classes with distinct
exposure vectors: βj1 ̸= βj2. Then there exists at least one sector i∗such that
δi := βij1 −βij2 ̸= 0
for some i = i∗.
(30)
Define the difference in automation exposure between the two labor classes
∆(t) := γj1(t) −γj2(t) =
n
X
i=1
δiαi(t),
(31)
where δi = βij1 −βij2. Since both weight vectors sum to one, we have P
i δi = 0, implying
that ∆(t) is a linear combination of αi(t) with zero-sum weights. Because at least one δi ̸= 0
and the αi(t) are strictly increasing, nonlinear, and parametrically heterogeneous (i.e., with
differing κi and/or t∗
i ), the weighted sum ∆(t) cannot be constant in t unless all αi(t) are
identical, which they are not under the model’s assumptions. Therefore, ∆(t) ̸≡0. By
continuity of ∆(t), it follows that
∃t > 0 such that ∆(t) ̸= 0
⇒
γj1(t) ̸= γj2(t).
(32)
Hence, automation exposure is strictly asymmetric across labor classes with differing sectoral
exposure profiles.
Proposition 3.1 formalizes how differences in sectoral employment patterns (βj) translate
into heterogeneous exposure to automation (γj(t)) across labor classes. The exposure func-
tion γj(t) = P
i βijαi(t) is a convex combination of sectoral automation intensities αi(t), with
weights βij. These weights reflect how much of labor class j’s employment is concentrated
in sector i. Since the AGI intensity αi(t) evolves differently across sectors due to differing
adoption rates (κi) and timing (t∗
i ), classes with different sectoral exposure weights will ex-
perience different composite automation trajectories. For instance, consider two classes j1
and j2, with βi∗j1 > βi∗j2 for a sector i∗that automates early (small t∗
i∗) or rapidly (large κi∗).
Then, the exposure path γj1(t) will rise faster than γj2(t), implying earlier and more intense
displacement pressure on class j1. Economically, this result captures the essence of differen-
tial automation risk: labor classes more concentrated in early-adopting or fast-automating
sectors face quicker erosion of wage income. As γj(t) enters directly into each household’s
12
This preprint research paper has not been peer reviewed. Electronic copy available at: https://ssrn.com/abstract=5333703
Preprint not peer reviewed
budget constraint (cf. Equation (16)), it reduces the effective labor income (1 −γj(t))wj(t),
tightening consumption possibilities unless offset by redistribution. This creates urgency for
targeted fiscal responses like class-calibrated UBI.n The proposition thus provides the the-
oretical foundation for class-specific displacement dynamics and justifies asymmetric policy
design in response to technological diffusion.
Proposition 3.2 (VAT Capture of AGI-Attributed Output Value). Let sector i produce
output Qi(t) via the nested CES production function defined in equation (8), where AGI
labor enters as
Li,AGI(t) = ψiαi(t)K(AGI)
i
(t),
(33)
with ψi > 0 representing AGI productivity, αi(t) ∈(0, 1) the automation intensity, and
K(AGI)
i
(t) the AGI capital stock.
Assume a VAT is levied at rate θi ∈[0, 1] on the revenue share attributed to AGI labor
in sector i. If the marginal product of AGI labor is
MP AGI
i
(t) :=
∂Qi(t)
∂Li,AGI(t),
(34)
then total VAT revenue from AGI-enabled production in sector i is
VATi(t) = θi · Pi(t) · MP AGI
i
(t) · Li,AGI(t).
(35)
Setting θi = 1 implies full taxation of AGI labor’s marginal revenue product.
Proof. From equation (33), AGI labor in sector i is given by
Li,AGI(t) = ψiαi(t)K(AGI)
i
(t).
Assuming perfect competition and constant returns to scale, factor payments equal marginal
products times quantities. The marginal product of AGI labor is
MP AGI
i
(t) =
∂Qi(t)
∂Li,AGI(t),
and the corresponding income contribution of AGI labor is
RevAGI
i
(t) := Pi(t) · MP AGI
i
(t) · Li,AGI(t).
Imposing a VAT at rate θi on this AGI-attributed revenue yields:
VATi(t) = θi · RevAGI
i
(t) = θi · Pi(t) · MP AGI
i
(t) · Li,AGI(t),
as stated.
When θi = 1, the VAT captures the full marginal revenue product of AGI
labor.
13
This preprint research paper has not been peer reviewed. Electronic copy available at: https://ssrn.com/abstract=5333703
Preprint not peer reviewed
Proposition 3.2 establishes that a value-added tax (VAT) levied on the marginal revenue
product of AGI labor provides a tractable mechanism for capturing the automation-induced
surplus in production. By linking the tax base directly to AGI-attributed output, the for-
mulation ensures that as AGI substitutes for human labor, the fiscal base for redistribution
shifts endogenously from wage income to automation-derived value. Setting θi = 1 corre-
sponds to fully taxing the marginal revenue generated by AGI labor, thereby creating a
revenue-neutral foundation for financing universal basic income (UBI) transfers calibrated
to the scale of labor displacement. Importantly, because the tax targets output rather than
capital directly, it preserves investment incentives while maintaining the state’s capacity to
redistribute gains from automation as the labor share in production declines.
Proposition 3.3 (Targeted UBI Sustainability Under Budget Balance). Let UBI(t) denote
the per-capita transfer to displaced workers at time t, and let
Ndisp(t) =
X
j
Njγj(t)
(36)
denote the total number of displaced individuals. Let aggregate VAT revenue collected from
all sectors be
RVAT(t) =
n
X
i=1
VATi(t).
(37)
Then, under a balanced-budget policy in which all VAT revenue is redistributed exclusively
to displaced individuals, a targeted UBI is fiscally sustainable if and only if
UBI(t) = RVAT(t)
Ndisp(t) .
(38)
Proof. The total fiscal expenditure on UBI at time t is given by
EUBI(t) = Ndisp(t) · UBI(t),
(39)
where Ndisp(t) is the number of displaced individuals as defined in (36). Total VAT revenue
is defined as
RVAT(t) =
n
X
i=1
VATi(t),
(40)
where each VATi(t) denotes the revenue collected from AGI-enabled production in sector i.
Under a pointwise balanced-budget condition—that is, assuming the government balances
its budget at each point in time—we require
EUBI(t) = RVAT(t).
(41)
Substituting (39) into (41) yields
Ndisp(t) · UBI(t) = RVAT(t).
14
This preprint research paper has not been peer reviewed. Electronic copy available at: https://ssrn.com/abstract=5333703
Preprint not peer reviewed
Solving for UBI(t), we obtain
UBI(t) = RVAT(t)
Ndisp(t) .
(42)
This identity must hold for UBI to be exactly funded by VAT revenue. Conversely, if (38)
holds, then government expenditures exactly match revenues by construction. Therefore, the
condition is both necessary and sufficient for targeted UBI sustainability under a balanced-
budget policy.
Equation (38) formalizes the direct relationship between automation-induced fiscal sur-
plus and income replacement for technologically displaced workers. Crucially, the formula
indicates that per-capita UBI is a function not of the total population, but of the mass of
displaced labor Ndisp(t) and the size of the VAT base. This formulation ensures that the
economic gains from AGI deployment are recycled specifically to those bearing the costs of
displacement, thereby operationalizing a principle of targeted fiscal fairness without diluting
transfers across the broader population.
However, the structure also introduces an inherent trade-off: as displacement intensi-
fies—reflected in higher γj(t) across labor classes—the denominator in (38) increases, which
may depress per-worker UBI unless VAT revenues grow proportionally. This dynamic creates
potential volatility in redistribution outcomes and underscores the importance of adaptive
policy design.
In this context, several fiscal mechanisms become critical. First, VAT rates may need
to adjust endogenously over time, or the AGI-related tax base may need to be broadened
to maintain revenue adequacy. Second, fiscal buffers or countercyclical borrowing can help
smooth UBI payments across periods of rapid automation. Third, integrating automation
monitoring into fiscal planning processes ensures that real-time labor market data informs
tax-and-transfer calibration.
Collectively, these measures support the sustainability and
equity of the redistribution scheme as the economy transitions toward higher levels of au-
tomation.
Equation (38) offers a transparent and operationalizable mechanism for aligning the fiscal
incidence of automation with its distributive consequences, thereby providing a foundation
for automation-responsive social policy.
Proposition 3.4 (Declining Labor Income Share with AGI Substitution). Consider a gen-
eral equilibrium economy in which each sector i produces output Qi(t) according to the
nested CES function in Equation (8), combining effective human labor Li,h(t) and AGI labor
Li,AGI(t) with elasticity parameter ρ1 < 1. Suppose:
(i) AGI intensity satisfies αi(t) →1 as t →∞for all i (Equation (3)),
(ii) Human labor supply declines such that Li,h(t) →0 as t →∞(Equations (5), (6)),
(iii) Wages are determined by the marginal product of human labor (Equation (18)),
(iv) Aggregate output P
i Qi(t) remains strictly positive as t →∞.
15
This preprint research paper has not been peer reviewed. Electronic copy available at: https://ssrn.com/abstract=5333703
Preprint not peer reviewed
Then, the aggregate labor income share converges to zero
lim
t→∞
P
i
∂Qi(t)
∂Li,h(t) · Li,h(t)
P
i Qi(t)
= 0.
(43)
Proof. From Equation (5), effective human labor in sector i is given by
Lij(t) = βijNj(1 −γj(t)).
As AGI intensity αi(t) →1, automation exposure γj(t) →1, implying
Li,h(t) =
X
j
Lij(t) →0
(by Equation (6))
AGI labor evolves as
Li,AGI(t) = ψiαi(t)K(AGI)
i
(t),
and since αi(t) →1 and K(AGI)
i
(t) grows over time (per Equation (13)), the AGI labor
input becomes large. In the CES aggregator with ρ1 < 1, human and AGI labor are gross
substitutes. As Li,h(t) →0 and Li,AGI(t) increases, the marginal product of human labor
∂Qi(t)
∂Li,h(t) →0. Therefore, sectoral labor income satisfies
∂Qi(t)
∂Li,h(t) · Li,h(t) →0.
Summing across all sectors,
X
i
∂Qi(t)
∂Li,h(t) · Li,h(t) →0.
By assumption, total output P
i Qi(t) remains bounded away from zero. Hence, the aggre-
gate labor income share converges to zero
lim
t→∞
P
i
∂Qi(t)
∂Li,h(t) · Li,h(t)
P
i Qi(t)
= 0.
Proposition 3.4 formalizes a central implication of AGI-driven automation: as human
labor becomes increasingly displaced (cf.
Equation (25)), the share of economic output
accruing to workers asymptotically vanishes. The nested CES production structure in Equa-
tion (8) implies that even when human and AGI labor are imperfect substitutes (ρ1 < 1), the
effective abundance and rising productivity of AGI labor enable it to substitute for human
input at scale. As AGI intensity rises and human labor supply contracts, the marginal prod-
uct of human labor declines (Equation (18)), driving wages toward zero despite continued or
even increasing aggregate output. In the long-run equilibrium, value-added becomes almost
entirely attributable to AGI-enabled capital services.
16
This preprint research paper has not been peer reviewed. Electronic copy available at: https://ssrn.com/abstract=5333703
Preprint not peer reviewed
In the model framework, households do not directly own physical or AGI-specific capi-
tal (see discussion following Equation (17)), implying that absent redistribution, displaced
workers face a collapse in income and consumption. This result underscores both the fiscal
and normative rationale for automation-contingent transfers, such as the VAT-financed, tar-
geted universal basic income (UBI) mechanism introduced in Equations (19)–(22). As the
labor income share declines, a robust and responsive social safety net becomes essential to
sustain household welfare and economic stability in the face of structural transformation.
Proposition 3.5 (S-Curve Dynamics of Labor Displacement). Let AGI intensity in sector
i evolve according to the logistic diffusion function (Equation (3))
αi(t) =
1
1 + e−κi(t−t∗
i ),
κi > 0.
(44)
Define the automation exposure of labor class j as
γj(t) =
n
X
i=1
βijαi(t),
with
X
i
βij = 1,
(45)
and let the effective labor supplied by class j to sector i be
Lij(t) = βijNj(1 −γj(t)).
(46)
Then the function Lij(t) exhibits an S-shaped trajectory over time, with curvature
d2Lij(t)
dt2





> 0
for t < t∗,
= 0
for t = t∗,
< 0
for t > t∗,
(47)
where t∗= P
i βijt∗
i is the weighted average of sectoral inflection points.
Proof. Differentiate Lij(t) = βijNj(1 −γj(t)) with respect to t:
dLij(t)
dt
= −βijNj
dγj(t)
dt
.
(48)
From the definition of γj(t), we have
dγj(t)
dt
=
n
X
k=1
βkj
dαk(t)
dt
,
(49)
and using the derivative of the logistic function
dαk(t)
dt
= κkαk(t)(1 −αk(t)).
(50)
17
This preprint research paper has not been peer reviewed. Electronic copy available at: https://ssrn.com/abstract=5333703
Preprint not peer reviewed
Thus,
dLij(t)
dt
= −βijNj
n
X
k=1
βkjκkαk(t)(1 −αk(t)).
(51)
Differentiating again
d2Lij(t)
dt2
= −βijNj
n
X
k=1
βkj
d
dt [κkαk(t)(1 −αk(t))]
(52)
= −βijNj
n
X
k=1
βkjκ2
kαk(t)(1 −αk(t))(1 −2αk(t)).
(53)
Note that each term in the sum changes sign at αk(t) = 0.5, i.e., t = t∗
k, because
1 −2αk(t)





> 0
for t < t∗
k,
= 0
for t = t∗
k,
< 0
for t > t∗
k.
The overall sign of d2Lij(t)
dt2
therefore switches from positive to negative as the weighted average
of inflection points is crossed. Since γj(t) is a convex combination of sigmoid functions αk(t),
its curvature changes at
t∗:=
X
k
βkjt∗
k,
yielding a global S-shaped transition for Lij(t), with inflection at t∗, as claimed.
Proposition 3.5 establishes that AGI-driven labor displacement follows a smooth, nonlin-
ear trajectory characterized by S-curve dynamics. This adjustment path accelerates around
a critical threshold t∗, defined as the weighted average of sector-specific automation inflec-
tion points. The S-shape arises from the logistic diffusion of AGI technologies within each
sector (Equation (3)), which induces a convex–concave transition in automation exposure
γj(t), and hence in the effective labor supply Lij(t).
Economically, this implies that labor displacement unfolds unevenly over time.
The
process begins gradually, reflecting low initial AGI intensity, then accelerates sharply as
automation technologies mature and diffuse, before tapering off as most displaceable human
labor is substituted or marginalized. The turning point t∗thus marks the period of maximum
displacement intensity.
From a policy perspective, this curvature has direct implications for the timing and
design of automation-responsive redistribution. The inflection point t∗offers a forecast of
when redistributive needs—such as those addressed through targeted UBI transfers (Equa-
tion (22))—will reach their peak. Anticipating the nonlinear dynamics of Lij(t) is therefore
essential for constructing fiscally adequate and temporally responsive interventions to miti-
gate sudden declines in labor income and to support macroeconomic stability during periods
of accelerated automation.
18
This preprint research paper has not been peer reviewed. Electronic copy available at: https://ssrn.com/abstract=5333703
Preprint not peer reviewed
Proposition 3.6 (Rising Income Inequality Under AGI Capital Concentration). Let to-
tal income in the economy be divided between labor income and AGI-based capital income.
Assume that:
Labor income accrues to the broad working population and is given by
YL(t) =
X
i
wi(t)Li,h(t),
(54)
where wi(t) denotes the wage in sector i, and Li,h(t) is the effective supply of human labor
(as in Equations (18) and (6)).
AGI capital income accrues to a fixed, small share of agents and is given by
YK(t) =
X
i
rAGI(t)K(AGI)
i
(t),
(55)
where K(AGI)
i
(t) evolves according to Equation (13).
Then, in the absence of redistribution and under AGI diffusion αi(t) →1, the Gini
coefficient G(t) rises over time:
dG(t)
dt
> 0
as t →∞.
(56)
Proof. Total income in the economy at time t is:
Y (t) = YL(t) + YK(t).
(57)
From Equation (6), human labor in sector i is:
Li,h(t) =
X
j
βijNj(1 −γj(t)),
(58)
where γj(t) = P
i βijαi(t) increases as AGI intensity αi(t) →1 (see Equation (3)). Therefore:
dγj(t)
dt
> 0
⇒
dLi,h(t)
dt
< 0.
(59)
Assuming wages wi(t) are non-increasing or weakly elastic to labor input, labor income
declines:
dYL(t)
dt
=
X
i
dwi(t)
dt
Li,h(t) + wi(t)dLi,h(t)
dt

< 0.
(60)
Simultaneously, AGI capital income increases with rising capital stock:
dYK(t)
dt
> 0,
(61)
given sustained growth in K(AGI)
i
(t) and persistently positive returns rAGI(t).
Since capital ownership is concentrated while labor income is broadly distributed, the
income distribution becomes increasingly unequal. The Gini coefficient G(t) therefore rises
over time, as formalized in Equation (56).
19
This preprint research paper has not been peer reviewed. Electronic copy available at: https://ssrn.com/abstract=5333703
Preprint not peer reviewed
Equation (56) formalizes a structural dynamic: labor income YL(t) declines (per Equa-
tion (60)) while capital income YK(t) increases (see Equation (61)), leading to growing
income concentration. As AGI labor substitutes for human labor and firms accumulate AGI
capital (via Equation (13)), income increasingly accrues to capital holders. In the absence
of redistribution—such as targeted UBI and AGI-linked VAT as defined in Equations (22)
and (20)—this transition leads to a mechanically increasing Gini coefficient. The result un-
derscores the need for automation-contingent fiscal tools to preserve distributional equity in
post-labor economies.
Proposition 3.7 (Targeted UBI Fully Offsets AGI-Induced Labor Income Loss). Let wj(t)
denote the wage of labor class j, as defined in Equation (18), and let γj(t) ∈[0, 1] be the
automation exposure of class j, from Equation (25). Define per-capita labor income for class
j as
yj(t) = (1 −γj(t))wj(t).
(62)
Suppose the government collects total VAT revenue,
RVAT(t) =
X
i
Ti(t),
(63)
as defined in Equation (20), and distributes it equally among the displaced population via a
targeted UBI
UBI(t) =
(RVAT(t)
Ndisp(t),
if Ndisp(t) > 0,
0,
otherwise,
(64)
where Ndisp(t) = P
j Njγj(t) is the total displaced population (Equation (21)).
Then, a sufficient condition for UBI to fully offset AGI-induced labor income losses for
all displaced individuals is
UBI(t) ≥max
j
[γj(t)wj(t)] .
(65)
Under this condition, the total income of every class j satisfies
yj(t) + UBI(t) ≥wj(t),
(66)
i.e., no group experiences a reduction in income relative to their pre-automation wage base-
line.
Proof. Per Equation (62), wage income for class j at time t is given by
yj(t) = (1 −γj(t))wj(t).
The corresponding income loss due to automation is
∆j(t) = wj(t) −yj(t) = γj(t)wj(t).
20
This preprint research paper has not been peer reviewed. Electronic copy available at: https://ssrn.com/abstract=5333703
Preprint not peer reviewed
Suppose UBI satisfies the sufficient condition in Equation (65)
UBI(t) ≥max
j
∆j(t).
Then, for every class j,
UBI(t) ≥∆j(t).
Adding UBI to the remaining labor income yields
yj(t) + UBI(t) ≥yj(t) + ∆j(t) = wj(t),
which confirms that total income does not fall below the pre-displacement level, proving
Equation (66).
Proposition 3.7 demonstrates that a targeted universal basic income (UBI) can fully offset
the loss in labor income caused by AGI-driven automation, provided the per-capita transfer
is at least as large as the maximum wage loss across labor classes. This condition defines a
sufficient upper bound for preserving individual income levels: if met, every worker’s total
income (wages plus UBI) remains no lower than it would have been absent automation.
Economically, this result offers a practical benchmark for designing redistribution poli-
cies aimed at neutralizing the distributive effects of technological displacement. However,
its applicability is subject to two key limitations. First, the condition may not be fiscally
feasible in scenarios of widespread displacement. If automation exposure γj(t) is high across
most classes, or if AGI deployment does not generate sufficiently large taxable surpluses,
then the VAT revenue required to fund the prescribed UBI could exceed available resources.
Second, the condition ensures income replacement but not welfare equivalence. Many work-
ers may value employment for reasons beyond income—such as social identity, structure, or
purpose—and the loss of work may impose utility costs that income alone cannot offset. Ad-
dressing these nonpecuniary losses would require complementary measures beyond monetary
compensation.
Despite these caveats, Proposition 3.7 provides an analytically grounded threshold for
automation-contingent redistribution.
It underscores how the sustainability of fiscal re-
sponses to technological change hinges on the evolving interplay between labor displacement,
AGI productivity, and the capacity of tax instruments to capture the surplus generated by
automation.
Proposition 3.8 (AGI-VAT Revenue as a Substitute for Displaced Payroll Base). Let wj(t)
denote the wage of labor class j at time t, and γj(t) denote the automation exposure of class
j, as defined in Equation (25). Then, the total displaced labor income at time t is
W loss(t) =
X
j
Njγj(t)wj(t),
(67)
where Nj is the size of class j. Let AGI labor input in sector i be defined via Equation (7)
Li,AGI(t) = ψiαi(t)K(AGI)
i
(t),
(68)
21
This preprint research paper has not been peer reviewed. Electronic copy available at: https://ssrn.com/abstract=5333703
Preprint not peer reviewed
and define the marginal product of AGI labor as
MP AGI
i
(t) =
∂Qi(t)
∂Li,AGI(t),
(69)
where Qi(t) is sectoral output from Equation (8).
Then the AGI-attributed output in sector i is
QAGI
i
(t) = MP AGI
i
(t) · Li,AGI(t).
(70)
If the government levies a value-added tax (VAT) on AGI output at rate θi ∈[0, 1] in
each sector, then VAT revenue from sector i is
Ti(t) = θi · QAGI
i
(t).
(71)
Total VAT revenue from AGI at time t is:
T(t) =
X
i
Ti(t).
(72)
Then a sufficient condition for the targeted UBI to fully offset displaced labor income is
T(t) ≥W loss(t),
(73)
i.e., the VAT revenue from AGI labor equals or exceeds the total displaced wage bill.
Proof. The displaced wage income is given by Equation (67), which aggregates lost income
from each labor class weighted by its exposure γj(t). The AGI labor in sector i, given by
Equation (68), contributes to output with marginal product MP AGI
i
(t) per Equation (69),
yielding the AGI-attributed component of production in Equation (70)
QAGI
i
(t) = MP AGI
i
(t) · Li,AGI(t).
Applying a VAT at rate θi to this quantity gives sectoral revenue in Equation (71), and
aggregating across sectors produces the total VAT revenue in Equation (72). If this total
revenue satisfies the condition in Equation (73), then
T(t) =
X
i
θiMP AGI
i
(t) · Li,AGI(t) ≥
X
j
Njγj(t)wj(t) = W loss(t),
which implies that the AGI-VAT system generates enough fiscal capacity to fully compensate
for displaced labor income.
Proposition 3.8 shows that under appropriate tax design, the government can fully reclaim
the labor value lost to automation by taxing AGI’s marginal contribution to production. This
acts as a fiscal mirror of the declining payroll base. This condition underpins Equation (23),
which guarantees that the targeted UBI is adequate to replace average foregone labor income.
Unlike payroll taxes, which shrink as labor is displaced, AGI-linked VAT dynamically scales
with automation intensity (via αi(t)) and capital deepening (K(AGI)
i
(t)), maintaining fiscal
adequacy in the face of structural change.
22
This preprint research paper has not been peer reviewed. Electronic copy available at: https://ssrn.com/abstract=5333703
Preprint not peer reviewed
Proposition 3.9 (Existence of a Post-Labor Balanced Growth Path). Consider an economy
where human labor is fully displaced, i.e., Li,h(t) = 0 for all sectors i and times t, and AGI
labor is fully active such that αi(t) →1 as t →∞, per Equation (3). Let AGI labor be
defined by
Li,AGI(t) = ψiαi(t)K(AGI)
i
(t).
(74)
Assume sectoral output is given by a CES function over AGI capital and AGI-derived labor
Qi(t) = Ai
h
ai

K(AGI)
i
(t)
ρ
+ (1 −ai) (Li,AGI(t))ρi 1
ρ ,
(75)
with ai ∈(0, 1), ρ < 1, and total factor productivity Ai > 0. Let capital accumulate according
to
˙K(AGI)
i
(t) = siQi(t) −δiK(AGI)
i
(t),
(76)
where si ∈(0, 1) is the savings rate and δi > 0 is the depreciation rate.
Then there exists a constant growth rate gi > 0 such that, along the balanced growth path
(BGP),
˙Qi(t)
Qi(t) =
˙K(AGI)
i
(t)
K(AGI)
i
(t)
= gi,
(77)
i.e., both capital and output grow at the same exponential rate in the post-labor economy.
Proof. As αi(t) →1, Equation (74) implies
Li,AGI(t) →ψiK(AGI)
i
(t).
(78)
Substitute this into Equation (75)
Qi(t) = Ai
h
ai

K(AGI)
i
(t)
ρ
+ (1 −ai)

ψiK(AGI)
i
(t)
ρi 1
ρ
= Ai [ai + (1 −ai)ψρ
i ]
1
ρ K(AGI)
i
(t).
(79)
Define the constant
Ξi := Ai [ai + (1 −ai)ψρ
i ]
1
ρ > 0,
(80)
so that Equation (79) becomes
Qi(t) = ΞiK(AGI)
i
(t).
(81)
Substitute into the capital accumulation equation (76):
˙K(AGI)
i
(t) = siΞiK(AGI)
i
(t) −δiK(AGI)
i
(t)
= (siΞi −δi) K(AGI)
i
(t).
(82)
Define the net growth rate
gi := siΞi −δi.
(83)
If gi > 0, then K(AGI)
i
(t) = Ki0egit and, from Equation (81), Qi(t) = ΞiKi0egit. Hence,
Equation (77) follows.
23
This preprint research paper has not been peer reviewed. Electronic copy available at: https://ssrn.com/abstract=5333703
Preprint not peer reviewed
Proposition 3.9 establishes that even in the complete absence of human labor, the econ-
omy can sustain a balanced growth path. AGI labor services, as defined in Equation (74),
become fully capital-embodied, and the CES production function reduces to a linear form
in AGI capital (Equation (81)). The endogenous growth rate gi, derived in Equation (83),
depends solely on productivity-adjusted savings and depreciation, decoupling growth from
demographic or labor-market constraints. This result validates the internal consistency of
the model’s post-labor regime: the economy does not collapse with labor displacement, but
instead converges to a purely capital-driven path of exponential growth.
Proposition 3.10 (Transitional Disequilibrium Due to AGI Capital Adjustment Costs).
Consider an economy where AGI-specific capital K(AGI)
i
(t) accumulates via investment I(AGI)
i
(t),
subject to convex adjustment costs given by
ϕi

I(AGI)
i
(t), K(AGI)
i
(t)

= ϑ
2
I(AGI)
i
(t)
K(AGI)
i
(t)
!2
K(AGI)
i
(t),
(84)
as defined in Equation (14). Let the aggregate resource constraint be
X
i
Qi(t) =
X
j
NjCj(t) +
X
i
I(AGI)
i
(t) +
X
i
ϕi

I(AGI)
i
(t), K(AGI)
i
(t)

,
(85)
as in Equation (??), and assume output Qi(t) is increasing in automation intensity αi(t)
(see Equation (3)).
Then, during the transition to AGI-intensive production, there exists a time interval over
which
d
dt
 X
j
NjCj(t)
!
< 0,
(86)
i.e., aggregate consumption declines temporarily—even as output rises—due to reallocation
of resources toward costly AGI capital adjustment.
Proof. From the resource constraint (85), define aggregate consumption
Cagg(t) :=
X
j
NjCj(t) =
X
i
Qi(t) −
X
i
h
I(AGI)
i
(t) + ϕi

I(AGI)
i
(t), K(AGI)
i
(t)
i
.
(87)
Differentiating both sides with respect to time
dCagg(t)
dt
=
X
i
dQi(t)
dt
−
X
i
"
dI(AGI)
i
(t)
dt
+ d
dtϕi

I(AGI)
i
(t), K(AGI)
i
(t)
\#
.
(88)
By the chain rule, the time derivative of the adjustment cost function is
d
dtϕi = ∂ϕi
∂Ii
· ˙I(AGI)
i
(t) + ∂ϕi
∂Ki
· ˙K(AGI)
i
(t),
(89)
24
This preprint research paper has not been peer reviewed. Electronic copy available at: https://ssrn.com/abstract=5333703
Preprint 
\end{document}
